\section{Introduction}
\label{sec:intro}

Multi-step LLM workflows---where each inference step feeds the next---have emerged as a
standard architecture for deploying language models in real applications. Frameworks such
as LangChain \cite{chase2022}, LangGraph \cite{langgraph2024}, AutoGen \cite{wu2023},
and PromptFlow \cite{microsoft2023} have made such workflows easy to build. Yet they
remain difficult to debug, resume, and improve systematically: state is carried forward
through conversation histories, mutable objects, and framework-specific control logic,
so an intermediate checkpoint at step $k$ may silently depend on earlier context not
captured in the checkpoint itself. Failure recovery becomes brittle, replay hard to
trust, and debugging requires reconstructing large portions of prior execution.

Tooling can add persistence and tracing, but the underlying limitation remains: the
plan, its intermediate semantics, and the conditions for valid continuation are scattered
across code and prompts rather than expressed in a single interpretable form.

\textbf{What is missing is a reasoning medium} --- a language in which multi-step
reasoning can be explicitly represented, inspected, verified, executed, and refined.
A viable reasoning medium must be simultaneously human-readable, machine-executable,
and structurally enforceable.

\textbf{NormCode is that medium.} NormCode \cite{anon2025} is a semi-formal planning
language for LLM workflows satisfying six properties that together constitute a working
reasoning medium: Interpretable, Enforceable, Composable, Locally Addressable,
Generalizable, and Portable (Sect.~\ref{sec:normcode}). The scope rule is the
structural mechanism that guarantees every checkpoint is a \textbf{self-contained case}
with no implicit context dependency (Sect.~\ref{sec:normcode}).

\textbf{NormCode Canvas (v1.1.3)} is the deployed system that exposes both levels.
Three practical properties follow structurally from the reasoning medium:
constant-effort failure localization~(C1), pre-execution structural review~(C2), and
scope-bounded selective re-execution~(C3) ---
demonstrated in Sect.~\ref{sec:canvas}--\ref{sec:casestudies}.

\textbf{Relation to the companion language paper \cite{anon2025}.} The companion paper
is a language specification: it defines NormCode's grammar, symbol set, type system,
and scope rule, and gives a formal account of the execution architecture. That paper
answers \textit{what NormCode is}. This paper answers \textit{what NormCode makes
possible}. In demonstration of its applicability, we present four production-ready
workflows ranging from PPT generation and code assistance to self-hosted NormCode
compilation.

\textbf{Contributions:}
(1)~Two-level CBR realized by a reasoning medium (Sect.~\ref{sec:level1}--\ref{sec:level2}):
Level~1 cases as suspended runtimes; Level~2 cases as runtime definitions; the compilation
pipeline as distillation; self-hosting as recursive case-based learning.
(2)~The six properties of a reasoning medium (Sect.~\ref{sec:normcode}): Interpretable,
Enforceable, Composable, Locally Addressable, Generalizable, Portable.
(3)~Three measurable properties of the deployed system
(Sect.~\ref{sec:canvas}--\ref{sec:casestudies}): failure localization (C1),
pre-execution review (C2), and selective re-execution (C3).
(4)~Four production-ready workflows (Sect.~\ref{sec:casestudies}): PPT generation,
code assistance, self-hosted compilation, and interactive canvas control.
